\section{What carries forward}
\label{sec:chapter-synthesis}

The chapter leaves later models with a small set of results that can be reused
without rebuilding their probabilistic foundations.  For a discrete binary
lineage, survival is governed by the exact recursion
\(
  S_{n+1}=2pS_n-pS_n^2
\).  It separates the three growth regimes, gives the critical tail
\(S_n\sim2/n\), and, below criticality, defines the survival amplitude \(A(p)\).
Conditioning on survival then produces a Yaglom limit with mean \(1/A(p)\), while
a finite founding cohort changes the chance of reaching late time but not this
limiting conditional scale.

The constant \(A(p)\) is available at three levels of resolution.  Its product
and series representations are exact; the bounds and the interval
\eqref{eq:A-certified-interval} provide certified finite computations; and the
asymptotic \(A(p)\sim2(1-2p)\) explains the divergence near criticality.  The
continuous-time constant
\(
  A_{\mathrm c}(p)=(1-2p)/(1-p)
\)
is not interchangeable with \(A(p)\): their difference records the available
surviving states under asynchronous and generational updates.  Later applications
must therefore choose the clock before importing a conditional population scale.

Catastrophe adds a separate selection mechanism.  The weighted recursion
\eqref{eq:cat-weighted-recursion} shows that conditioning on avoidance of rupture
reweights entire population paths, not only their terminal counts.  A
quasi-stationary result for ordinary branching consequently cannot be transferred
to a killed or rupturing process without defining its transient kernel and
conditioning event.  This point becomes more important, not less, in two-type
models, where the catastrophe rate may depend on the composition as well as the
total census.

For transfer between compartments, the bivariate PGF keeps the interior and
exterior counts distinct.  Pure absorption and absorption--death give binomial
and trinomial laws, while absorption--birth--death has the global integral
solution \eqref{eq:abd-integral-solution} and, away from critical and resonant
parameters, the characteristic form \eqref{eq:Gabd}.  In a multi-compartment
extension, the same construction can be applied after a rupture by replacing the
single exterior cohort with the released burst-size distribution and assigning
one absorption channel to each receiving compartment.  Multivariate PGFs then
record the resulting joint counts; no new probability principle is required,
although the state space and characteristic geometry become larger.

These results also mark the boundary between model and application.  The
mathematics states what follows from independent reproduction, constant
state-dependent hazards and the chosen event ordering.  Whether those assumptions
are adequate for a particular population, compartment or rupture mechanism must
be decided in the relevant application chapter.  The quantities to test are now
specific: early extinction probabilities, the survival-conditioned scale,
catastrophe-time and burst-size laws, absorption times, and the sensitivity of
each to distance from criticality.
